\documentclass[11pt]{article}

\usepackage[a4paper,margin=2.5cm]{geometry}
\usepackage{amsmath}
\usepackage{hyperref}

\title{LLP Acceptance Lab: un estudio fenomenologico toy de particulas de larga vida}
\author{}
\date{}

\begin{document}

\maketitle

\section{Introduccion}

% TODO: Add real references and a short introduction to long-lived particles.

\section{Cinematica relativista}

For a particle with total energy $E$ and mass $m$ in natural units,
\begin{equation}
  \gamma = \frac{E}{m},
  \qquad
  \beta\gamma = \frac{p}{m}
  = \frac{\sqrt{E^2 - m^2}}{m}.
\end{equation}

\section{Probabilidad de decaimiento espacial}

The mean decay length in the lab frame is
\begin{equation}
  L_{\mathrm{lab}} = \beta\gamma c\tau.
\end{equation}

The probability to decay between two path lengths $L_1$ and $L_2$ is
\begin{equation}
  P(L_1 < L < L_2)
  = \exp\left(-\frac{L_1}{L_{\mathrm{lab}}}\right)
  - \exp\left(-\frac{L_2}{L_{\mathrm{lab}}}\right).
\end{equation}

\section{Regiones detectoras idealizadas}

% TODO: Add toy regions in a later phase. Do not claim real detector dimensions.

\section{Distribuciones de boost}

% TODO: Add boost-averaged acceptance.

\section{Geometria cilindrica toy}

% TODO: Add cylindrical toy geometry.

\section{Aplicacion toy a una ALP}

% TODO: Add toy ALP-like photon-coupling lifetime.

\section{Conclusiones}

% TODO: Summarize lessons and limitations after the implementation is complete.

\bibliographystyle{unsrt}
\bibliography{references}

\end{document}
